\documentclass{article}
\usepackage[T1]{fontenc}
\usepackage{lmodern}
\usepackage{graphicx}
\usepackage{amsmath,amssymb}
\usepackage[a4paper,margin=2cm]{geometry}
\usepackage[absolute,overlay]{textpos}
\setlength{\TPHorizModule}{1mm}
\setlength{\TPVertModule}{1mm}
\usepackage[indonesian]{babel}
\usepackage{hyperref}
\setlength{\emergencystretch}{2em}
% unit: Halaman 10

\begin{document}

% === Halaman 10 (llm) ===

\section*{Algoritma \textit{Rijndael}}

\begin{itemize}
    \item Yang dijelaskan di sini Rijndael 128-bit (anjang kunci 128 bit)
    \item Tidak seperti \textit{DES} yang berorientasi bit, \textit{Rijndael} beroperasi dalam orientasi \textit{byte}.
    \item \textit{Rijndael} tidak menggunakan jaringan Feistel seperti cipher blok lainnya.
    \item \textit{Enciphering} dilakukan dalam sejumlah putaran (iterate cipher). Setiap putaran mengunakan kunci internal yang berbeda (disebut \textit{round key}).
    \item \textit{Enciphering} melibatkan operasi substitusi dan permutasi.
\end{itemize}

\end{document}
